% GENERATED FILE. Edit canonical lessons, then run npm run generate:books.
% canonical-source-hash: fnv1a64:17717a752838f9f1
% canonical-lessons: MR-R25-twelve-verbs, MR-R25-dative-frame, MR-R25-recognition-script, MR-R25-recognition-words

\chapter{The Twelve, Further Apart}
\label{ch:mr-twelve-verbs-r2}
\hlchaptermodality{25}

\begin{chapteropening}
\textbf{By the end of this chapter:} \emph{I can produce all twelve of the new verbs at a wider spacing, and say what the mala frame does in each of the four sentences built on it.}

Complete MR-R25-dative-frame and produce feeling, sleepiness, and getting on one dative frame, then say word for word what mala is doing in each.
\end{chapteropening}

\section[bārā kriyāpadaṁ]{The twelve, further apart}

\label{lesson:MR-R25-twelve-verbs}



\begin{quote}
The four you ask for: \textbf{\mr{देणे}, \mr{पिणे}, \mr{आणणे}, \mr{ठेवणे}}. Say them without looking anywhere.
\end{quote}

\subsection*{Your turn}
Then the four that pass between two people: \emph{sāṅgṇe}, \textbf{\mr{म्हणणे}, \mr{शिकणे}, \mr{मिळणे}}.

\begin{tcolorbox}[breakable,colback=teal!4,colframe=teal!35!black,title={Before you move on}]
Twelve verbs, three runs of four. Say the one you would use to accept an offer. (\textbf{Chālel.}) Say the one you would use to ask for water. (\textbf{Dyā}.) Then stop.
\end{tcolorbox}
\section[malā …]{One frame, four sentences}

\label{lesson:MR-R25-dative-frame}



\begin{quote}
\textbf{\mr{मला}} is where things arrive. Say the two oldest ones: \textbf{\mr{मला} \mr{माहीत} \mr{आहे}}, \textbf{\mr{मला} \mr{मराठी} \mr{आवडते}}.
\end{quote}

\subsection*{Your turn}
Say each one, then say the word-for-word English under it.

\begin{tcolorbox}[breakable,colback=teal!4,colframe=teal!35!black,title={Before you move on}]
What does \textbf{\mr{मला}} do in every one of these? (\textbf{It is where the knowing, the liking, the feeling and the getting arrive}.) And name the body word this book gave you for where feeling is said to sit. (\textbf{\mr{हृदय}}.) Then stop.
\end{tcolorbox}
\section[ha / o / chā]{Recognition at distance: two signs and one sound}

\label{lesson:MR-R25-recognition-script}



\begin{quote}
The first two shapes this book ever asked you to make were \textbf{\mr{ह}} and the \textbf{\mr{ो}} mark, and together they wrote \textbf{\mr{हो}}. Make them both from memory.
\end{quote}

\subsection*{Your turn}
Say \textbf{\mr{चालणे}} and \textbf{\mr{चार}} one after the other and hear the same leaning in both.

\begin{tcolorbox}[breakable,colback=teal!4,colframe=teal!35!black,title={Before you move on}]
Write \textbf{\mr{हो}} once without a model, say \textbf{\mr{चार}} with its leaning consonant, then stop.
\end{tcolorbox}
\section[āhe / mitra / kuṭumb / samajṇe]{Recognition at distance: five things from further back}

\label{lesson:MR-R25-recognition-words}



\begin{quote}
Count to five out loud — \emph{ek, don, tīn, chār, pāch} — and then put \textbf{\mr{आहे}} at the end of a sentence about yourself, where Marathi keeps it.
\end{quote}

\subsection*{Your turn}
Both were taken up from Sanskrit whole rather than worn down, and \textbf{\mr{कुटुंब}} takes a gender Hindi does not give it.

\begin{tcolorbox}[breakable,colback=teal!4,colframe=teal!35!black,title={Before you move on}]
Say that Marathi is understandable to you (\textbf{\mr{मला} \mr{मराठी} \mr{समजते}}) and then that you are learning it (\textbf{\mr{मी} \mr{मराठी} \mr{शिकतो} / \mr{शिकते}}). What did \textbf{\mr{मित्र}}'s root mean before it meant friend? (\textbf{To bind}.) Then stop.
\end{tcolorbox}
